\subsection{Scene II --- The Deformable Face}\label{subsec:scene_face}

The second scene, \texttt{Mario\_N64}, is the simplest in conception and the
most direct test of spatial precision. It recreates the grab-and-stretch
interactive face familiar from an early game console as a hand-tracking exercise:
a head mesh that the user pulls and distorts with the hands. \texttt{FaceDeformer}
works on a private copy of the mesh, so the shared asset is never altered, and
maintains for every vertex a rest position, a current position, a velocity, and a
per-frame target.

\begin{figure}[H]
	\centering
	\captionsetup{justification=centering}
	\includegraphics[width=0.9\linewidth, frame]{img/scene_face_deformer.png}
	\caption{The head mesh pulled out of shape
	         by two independent grabs.}
	\label{fig:face-deformer}
\end{figure}

A rendered mesh splits vertices at its seams, where a single physical point is
represented by several entries carrying different texture coordinates; deforming
such a mesh naively tears it along those seams. The component prevents this by
grouping vertices that share a position into welded sets at start-up, running the
physics on one representative per set, and copying the result to its duplicates.
The visible surface stays continuous under arbitrarily large pulls.

A grab begins when a hand closes near the surface. The nearest vertex becomes the
grab centre, and a smooth, Gaussian falloff assigns every other vertex a weight
that decays with distance from it, so the pull is concentrated at the grasp and
tapers across the surrounding region rather than dragging a single point. The two
hands grab independently, each with its own centre and weights, so the face can be
pulled in two places at once.

While a grab is held, the target positions of the affected vertices follow the
hand; between and after grabs, a mass-spring relation pulls every representative
vertex back toward its rest position, the restoring force damped exponentially so
the surface settles with an elastic wobble rather than snapping back. Normals and
bounds are recomputed each frame so the lighting tracks the deformation.

The user orbits the head through the same ring camera as the previous scene, and a
small interlock suppresses grabbing while the ring is actively turning, so rotating
the view cannot be mistaken for reaching into the mesh. Modest as it is, the scene
makes a point that bears directly on~\autoref{subsubsec:sota_gesture}: the recovered
pose is precise enough to pick out and drag an individual vertex continuously, in
three dimensions, with either hand.
